\subsection{Réversibilité de la constante rationnelle et changement de phase}

La constante rationnelle \(\kappa_{\mathrm{per}}\) détermine intégralement
le registre \(K\) lorsque la base, la longueur et l'origine
de lecture sont conservées. Sa fonction numérique et sa fonction structurelle sont reliées
par une inversion exacte.

\begin{proposition}[Récupération entière et transformation cyclique]
Soit \(I(K)=\sum_{j=1}^{12}K_j1000^{12-j}\). Alors
\[
 I(K)=(1000^{12}-1)\kappa_{\mathrm{per}}.
\]
Le développement entier de \(I(K)\) en douze blocs de base \(1000\)
récupère tous les \(K_j\), y compris les zéros initiaux de chaque bloc.
Si \(\rho(K)=(K_2,\ldots,K_{12},K_1)\), on a
\begin{equation}
 \kappa(\rho K)=1000\kappa(K)-K_1.
 \label{eq:k-rotacion-racional}
\end{equation}
\end{proposition}
\begin{proof}
La première identité inverse la définition de l'évaluation périodique.
La division euclidienne itérée par \(1000\), avec douze positions fixées,
récupère le vecteur. D'autre part,
\[
 I(\rho K)=1000I(K)-K_1(1000^{12}-1).
\]
La division par \(1000^{12}-1\) donne \eqref{eq:k-rotacion-racional}.
\end{proof}

La récupération compose donc \(\kappa\mapsto K\mapsto U\)
et les transformations \(D_3K,D_4K,Q\) déjà construites. L'évaluation
rationnelle conserve le registre signé et les différences cycliques qui
en dépendent. La profondeur totale et les opérateurs de transport d'une
histoire restent dans leurs coordonnées propres : l'inversion précédente
a pour domaine le registre entier dodécaphasique.

Écrivons \(\kappa_j=\kappa(\rho^jK)\), avec des indices modulo \(12\).
La même identité donne
\[
 K_{j+1}=1000\kappa_j-\kappa_{j+1},\qquad
 \sum_{j=0}^{11}\kappa_j=\frac{\sum_{j=1}^{12}K_j}{999}
 =\frac{6263}{999}.
\]
La phase agit de manière exacte sur la constante : sa valeur est fixe
dans la lecture canonique, tandis que ses rotations suivent une loi affine.
La somme de l'orbite est invariante sous le choix de l'origine.

L'incidence ultérieure possède également une expression dans ces
coordonnées. Le seuil cubique détermine
\[
 B_K=\{j+1:1000\kappa_j-\kappa_{j+1}\ge729\}
 =\{4,6,9,10\}.
\]
La tétrade se récupère à partir de l'orbite rationnelle parce que celle-ci
récupère d'abord les composantes entières. Le support signé associé
à \(\alpha\) utilise en outre ses registres régionaux et la normalisation
de frontière. La relation entre coordonnée et incidence s'exprime ainsi
au moyen d'opérations identifiables.

C'est une définition opérationnelle de la fonction de \(K\) : il fournit
une mise en coordonnées entières réversible du registre orienté, dont
l'évaluation périodique et les fonctions d'incidence participent aux
constructions ultérieures. La valeur, la transformation et le domaine
sont conservés conjointement.
